\usepackage[margin=2cm]{geometry}
\usepackage{amsmath, amsfonts, amssymb, xcolor, float, amsthm, mathtools, enumitem, cancel, soul, booktabs}
\usepackage[calcwidth]{titlesec}
\usepackage{ragged2e}

% \usepackage{graphicx}
% \graphicspath{ {./imgs/} }

\usepackage[hidelinks]{hyperref}
\hypersetup{
    colorlinks=true,
    citecolor=blue,
    filecolor=blue,
    linkcolor=blue,
    urlcolor=blue
}
\usepackage[nameinlink,capitalise,noabbrev]{cleveref}
\crefname{example}{Example}{Examples}
\crefname{equation}{Equation}{Equations}

\usepackage[T1]{fontenc}
\usepackage{minted}

\setminted{
    fontsize=\normalfont,
    style=emacs,
    linenos,
    breaklines,
    breakanywhere,
    autogobble,
    tabsize=4,
    frame=lines,
    framesep=1mm,
    numbersep=0.5em,      
}

% ── Python minted shortcuts ──────────────────────────────────
\newcommand{\py}[1]{\mintinline{python}|#1|}
\newcommand{\Rcode}[1]{\mintinline{R}|#1|}

% so verbatim tokenisation works correctly
\newminted[pyblock]{python}{}
\newminted[Rblock]{R}{}


% Optional: tweak line number look
\renewcommand\theFancyVerbLine{\sffamily\scriptsize\arabic{FancyVerbLine}}
\usepackage{fvextra}
\DefineVerbatimEnvironment{pyout}{Verbatim}{
    fontsize=\small,
    frame=lines,
    framesep=2mm
}

\newcommand{\R}{{\rm I\!R}}
\let\biconditional\leftrightarrow

%Formatting the sections
\renewcommand{\thesection}{\arabic{section}}
% \renewcommand{\thechapter}{\Roman{chapter}}
\setcounter{secnumdepth}{3}
\titleformat{\section}[hang]{\bfseries\Large}{\thesection}{1em}{}
\titleformat{\subsection}[hang]{\bfseries\filleft\large}{\itshape\thesubsection}{1em}{}[\vspace{-1.5ex}\rule{1.4\titlewidth}{1pt}]

\newcounter{example}[section] % resets each section (optional)
\renewcommand{\theexample}{\thesection.\arabic{example}} % e.g. 1.3.1 (optional)
\newcommand{\exampleparagraph}{%
  \refstepcounter{example}%
  \paragraph{Example \theexample}%
}

% Theoretical Statements Section (scrapped not useful)
% \titleclass{\statement}{straight}[\subsubsection]
% \newcounter{statement}
% \titleformat{\statement}[hang]{\bfseries}{\itshape\thestatement\quad--}{1em}{}[]
% \titlespacing{\statement}{1em}{1em}{1em}[1em]
% \renewcommand{\thestatement}{\arabic{statement}}

%Theoretical Statements Envs
%amsthm
\newtheoremstyle{break}%
    {}{}
    {}{\parindent}
    {\bfseries}{:}
    {\newline}{}
\theoremstyle{break}
\newtheorem{lem}{Lemma}
\newtheorem{cor}{Corollary}
\newtheorem{prop}{Proposition}
\newenvironment{lemma}[1]{\begin{lem}[#1]\leavevmode\vspace{-1.5\baselineskip}}{\end{lem}}
\newenvironment{corollary}[1]{\begin{cor}[#1]\leavevmode\vspace{-1.5\baselineskip}}{\end{cor}}
\newenvironment{proposition}[1]{\begin{prop}[#1]\leavevmode\vspace{-1.5\baselineskip}}{\end{prop}}

%For boxing environment
\usepackage{tcolorbox}
\tcbuselibrary{skins, breakable, most}
\newtcolorbox{quotebox}[1][]{%
  enhanced jigsaw,
  sharp corners,
  colback=white,
  borderline={1pt}{-2pt}{black},
  fontupper={\setlength{\parindent}{20pt}},
  #1
}
%Theorem box
\newtcolorbox{theorembox}[1]{
    breakable = true,
    title = #1,
    coltitle = black,
    colback = white,
    fonttitle = \bfseries,
    enhanced,
    sharp corners,
    boxrule = 0.5pt,
    colframe = black,
    attach boxed title to top left = {xshift = 15mm, yshift = -5.5mm},
    boxed title style={
        frame code={
            \path[draw = black, fill = black]
                ([yshift=-1mm,xshift=-1mm]frame.north west)
                arc[start angle=0,end angle=180,radius=1mm]
                ([yshift=-1mm,xshift=1mm]frame.north east)
                arc[start angle=180,end angle=0,radius=1mm];

            \path[draw = black, fill = white]
                ([xshift=-2mm]frame.north west) -- ([xshift=2mm]frame.north east)
                [rounded corners=1mm]-- ([xshift=1mm,yshift=-1mm]frame.north east)
                -- (frame.south east) -- (frame.south west)
                -- ([xshift=-1mm,yshift=-1mm]frame.north west)
                [sharp corners]-- cycle;
        },
        interior engine = empty,
        colback = white,
    },
    width = 0.95\linewidth,
    top = 10mm,
    left = 5mm,
    overlay = {

    }
}
% \newenvironment{theorembox}[1]
    % {\begin{center}\begin{theobox}{#1}}
    % {\end{theobox}\end{center}}

%Definition box
\newtcolorbox{definitionbox}{
    breakable = true,
    title = Definition,
    fonttitle = \bfseries,
    colback = white,
    width = 0.85\linewidth,
    coltitle = black,
    frame hidden,
    sharp corners,
    enhanced,
    attach boxed title to top left = {xshift = -0.5cm, yshift = -0.1mm},
    boxed title style = {
        colback = white,
        frame hidden,
    },
    overlay = {
        \node (A) at (title.south west) {};
        \node (B) at (title.south east) {};
        \node (C) at (frame.south west) {};
        \draw [semithick] (A) -- (B);
        \draw [semithick, double distance = 2pt] (C) -- (frame.north west);
    },
}
% \newenvironment{definitionbox}
    % {\begin{center}\begin{defbox}}
    % {\end{defbox}\end{center}}

\newtcolorbox{proofbox}{
    enhanced,
    breakable,
    colback=white,
    frame hidden,
    width=0.85\linewidth,
    left=2cm,                       % reserve space on the left
    overlay={
        % left vertical rule
        \draw[semithick] ([xshift=1.8cm]frame.north west) -- 
                        ([xshift=1.8cm]frame.south west);
        % "Proof" label at the top-left reserved area
        \node[anchor=north west,font=\itshape]
        at ([xshift=0.2cm,yshift=-0.2cm]frame.north west) {proof};
    },
}
% \newenvironment{proofbox}
    % {\begin{center}\begin{boxproof}}
    % {\end{boxproof}\end{center}}

\newtcolorbox{ProofIllustrationBox}{
    breakable = true,
    colback = white,
    coltitle = black,
    sharp corners,
    enhanced,
    %segmentation
    segmentation style = {
        solid,
        black,
    },
    %font
    fontlower = \sffamily\itshape\scriptsize,
    %format
    size = fbox,
    halign lower = center,
    halign upper = center,
}
\newtcolorbox{example}{
    breakable = true,
    title = Example,
    sidebyside,
    sidebyside align = top seam,
    enhanced,
    coltitle = black,
    colback = white,
    %Title
    attach boxed title to top left = {xshift = 1cm,},
    boxed title style = {
        colback = white,
        colframe = black,
        skin = enhancedfirst jigsaw,
        size = small,
        arc = 1mm,
        bottom = -1mm,
        boxrule = 0.8pt,
    },
    %Segmentation
    segmentation style = {
        thick,
        solid,
        shorten <= 10pt,
        shorten >= 10pt,
    },
    sharpish corners,
    colframe = black,
    boxrule = 0.8pt,
    width = 0.8\linewidth,
}
\newenvironment{examplebox}
    {\begin{center}\begin{example}}
    {\end{example}\end{center}}



%Tikz
\usepackage{tikz}
\usetikzlibrary{spy, intersections, calc, angles, quotes, fit, decorations.pathmorphing}
\tikzset{help lines/.style 2 args = {very thin, color = #1, step = #2}}
\tikzset{Hai lines/.style = {help lines = magenta}}

%For Pseudocode
\usepackage[ruled,titlenotnumbered]{algorithm2e}
\NoCaptionOfAlgo
\SetAlgoCaptionSeparator{}
\SetAlCapNameFnt{\large\scshape}
\DontPrintSemicolon
\SetAlgoLongEnd
\SetAlgoInsideSkip{smallskip}
\SetKwProg{Fn}{function}{:}{end function}
\SetKwComment{Comment}{$\triangleright$\ }{}
\usepackage{caption}
%Highlighting the pseudo-code in the algorithm environment using Tikz fit
\newcommand{\boxit}[3]{
    \tikz[remember picture,overlay] \node (A) {};\ignorespaces
    \tikz[remember picture,overlay]{
            \node[yshift=3pt,fill=#1,opacity=.25,
            fit={ ($(A)+(-0.1,0.3\baselineskip)$) ($(A)+(#2\textwidth,-{#3}\baselineskip - 0.25\baselineskip)$) } ] {};
        }\ignorespaces
}
\newcommand\mycommfont[1]{\footnotesize\ttfamily\itshape #1\unskip}
\SetCommentSty{mycommfont}
\setlength{\algomargin}{2em}
\setlength{\interspacetitleruled}{6pt}
\setlength\algoheightrule{1.5pt}


% \title{\texttt{Python} Notes}   
% \date{\today}
% \author{Tri Hai Huynh \\ \\
        % Undergraduate Student of The University \emph{of} Adelaide, \\ \\
        % Adelaide, South Australia}
